\scanpage{109}
% archive foliation: 108

la dérivation $\Theta_\delta$.

Inversement, partant de $a_1,a_2,a_3\in k$, \deleted{\unsure{satisfaisant}}
\inserted{tels que} le polynôme
\[
Z^3+a_1Z^2+a_2Z+a_3\in k[Z]
\]
soit sans racines multiples, i.e tel que
\[
(34)\qquad \underbrace{\Delta(a_1,a_2,a_3)}_{\text{discriminant}}\ne0,
\]
on vérifie que les formules précédentes définissent une courbe elliptique sur $k$
(qu'on peut \deleted{définir} \inserted{construire} comme le
$\operatorname{Proj}(\mathcal F_{\unsure{*|\infty}})$, où
$\mathcal F_\infty$ est défini par (31)), \deleted{\unsure{de $\mathcal F$ muni}}
\inserted{\unsure{tel que}} un \deleted{\unsure{morphisme}} \inserted{quadratique}
en $\mathbb P^1_k$ défini par
\[
(35)\quad \deleted{\unsure{\wp}}\qquad
\gamma:\diagram{\begin{tikzcd}[column sep=large,row sep=large]
E\arrow[r]\arrow[dr] & \mathbb P^1_k\ \deleted{\unsure{\wp}}\\
& E/\{\pm\mathrm{id}_E\}\arrow[u,"\simeq"']
\end{tikzcd}}
\]
ramifié en \deleted{des} \inserted{\unsure{l'ens.} $J$} trois points
\inserted{$z_i$} de $\mathbb P^1(k)\unsure{\setminus}\{\infty\}\inserted{{}=k}$
solutions de l'équation
\[
(36)\qquad Z^3+a_1Z^2+a_2Z+a_3=0,
\]
\deleted{\unsure{l'élément|l'origine}} \inserted{\unsure{l'élément|l'origine}} neutre de $E$ étant
au dessus de \deleted{\unsure{le point}} \inserted{\unsure{l'élément|le point}} $\infty$ de $\mathbb P^1_k$.
\deleted{qui est} Bien \deleted{\unsure{entendu}}\inserted{\unsure{sûr}}, l'inclusion $J\subset\mathbb P^1_k$ définit
un isom canonique
\[
(37)\qquad \Sigma_J\simeq\mathbb P^1_k,\qquad t\longmapsto\infty.
\]
